\documentclass{article}
\usepackage{amsmath,amsthm,amssymb}
\usepackage[margin=1in]{geometry}

\DeclareMathOperator{\area}{area}
\DeclareMathOperator{\lowest}{lowest}
\DeclareMathOperator{\sym}{sym}
\newtheorem{theorem}{Theorem}

\title{Computer-Assisted Proofs for a Rational \(q,t\)-Catalan Formula}
\author{}
\date{}

\begin{document}
\maketitle

\section{The theorem and the role of the computations}

This note explains the computer-assisted part of the proof in
Hawkes's paper \emph{A conjectured formula for the rational
\(q,t\)-Catalan polynomial}~\cite{Hawkes2024}.  The paper proposes a
symmetric-string formula for the rational \(q,t\)-Catalan polynomial.  The
formula remains conjectural in general, but the paper proves it through
bounded deficit in the family
\[
  r\equiv 1\pmod s.
\]
Write
\[
  s=\ell+1,
  \qquad
  r=(\ell+1)m+1,
  \qquad
  M_{\ell,m}=m\binom{\ell+1}{2}.
\]
The proved result is the following.

\begin{theorem}[Hawkes]
If \(r=(\ell+1)m+1\) with \(m\geq1\), then the conjectured
symmetric-string formula is correct in every deficit layer
\[
  0\leq d\leq20.
\]
\end{theorem}

The computation is not a search for supporting examples.  It verifies two
finite obligations that the symbolic proof isolates.  We call them
\emph{Computation A} and \emph{Computation B}:
\begin{itemize}
\item Computation A verifies the proposed formula directly in the finite base
      range.
\item Computation B verifies that the explicitly constructed canonical
      strings do not run past their prescribed endpoints.
\end{itemize}
The symbolic part of the proof then completes the canonical strings at one
critical length and transports the resulting decomposition to every larger
length.

More generally, the paper shows that for any fixed cutoff \(d^*\), producing
this kind of combinatorial proof for all layers \(d\leq d^*\) in the family
\(r\equiv1\pmod s\) is equivalent to a finite verification problem.  The
choice \(d^*=20\) is therefore not intrinsic to the method.  A larger cutoff
would require additional finite calculations, which have not been recorded
here.

\section{The proposed strings}

Let \(D_{\ell,m}^d\) denote the position-coordinate \(m\)-Dyck paths in
deficit layer \(d\), and let
\[
  C_{\ell,m}^d(q,t)
  =
  \sum_{x\in D_{\ell,m}^d}
  q^{\area(x)}t^{M_{\ell,m}-d-\area(x)}.
\]
The distinguished roots form a subset \(T_{\ell,m}^d\).  In the position
coordinates used by the code, a path is a sequence
\[
  x=(a_0,a_1,\ldots,a_\ell),
  \qquad
  a_0=0,
  \qquad
  0\leq a_{i+1}\leq a_i+m,
\]
and the roots in \(T_{\ell,m}^d\) are characterized by beginning
\[
  (0,0,\ldots).
\]

For integers \(a,b\), define the signed symmetric interval
\[
  \sym(a,b)
  =
  \frac{q^{b+1}t^a-q^a t^{b+1}}{q-t}.
\]
When \(a\leq b\), this is the ordinary positive string
\[
  q^a t^b+q^{a+1}t^{b-1}+\cdots+q^b t^a.
\]
The conjectured layer formula is
\[
  C_{\ell,m}^d(q,t)
  =
  \sum_{x\in T_{\ell,m}^d}
  \sym\bigl(\area(x),M_{\ell,m}-d-\area(x)\bigr).
  \tag{1}
\]
Thus every root prescribes a ``string holster'': one slot at every consecutive
area from the root area to its symmetric terminal area.

When \(d<(\ell-1)m\), the paper constructs the roots from fitting partitions.
If \(P_{\ell-1}^d\) is the set of partitions of \(d\) whose largest part is at
most \(\ell-1\), then maps \(g\) and \(f\) give a bijection
\[
  f\circ g:P_{\ell-1}^d\longrightarrow T_{\ell,m}^d.
\]
For \(m=1\), once \(\ell-1\geq d\), these fitting partitions are simply all
partitions of \(d\).

\section{Canonical initial strings}

For each fitting partition \(\lambda\), let
\[
  x_\lambda=f(g(\lambda)).
\]
The paper defines a local operation \(\operatorname{right}\).  Beginning at
\(x_\lambda\), repeatedly apply this operation until it is no longer
possible:
\[
  x_\lambda,
  \operatorname{right}(x_\lambda),
  \operatorname{right}^2(x_\lambda),
  \ldots,
  \lowest(x_\lambda).
\]
Each step raises area by one while preserving the deficit.  These are the
\emph{canonical initial strings}.

The following facts are proved symbolically in the paper:
\begin{itemize}
\item distinct canonical strings are disjoint;
\item the maps \(\operatorname{left}\) and \(\operatorname{right}\) are
      inverse where they are defined;
\item the canonical strings partition the subset
      \(C_{\ell,m}^d\subseteq D_{\ell,m}^d\) of connected paths.
\end{itemize}
The remaining paths
\[
  N_{\ell,m}^d=D_{\ell,m}^d\setminus C_{\ell,m}^d
\]
are available to complete the canonical strings.  The added tails need only
occupy consecutive area levels; they are not required to continue the
\(\operatorname{right}\) operation.

\section{The critical length}

Fix a cutoff \(d^*\).  For each \(m\), define
\[
  \ell^*(m,d^*)
  =
  \min\{\ell:(\ell-1)m>d^*\}.
\]
The proof has three stages:
\begin{enumerate}
\item verify the formula directly for \(\ell\leq\ell^*\);
\item prove that the canonical strings at \(\ell=\ell^*\) can be completed;
\item propagate that completed decomposition symbolically to every
      \(\ell>\ell^*\).
\end{enumerate}
For the theorem, \(d^*=20\).

\section{Computation A: the finite formula check}

Computation A is Lemma 3 of Section 9; it is called ``computation 2'' in the
paper's appendix.  For each required pair \((m,\ell)\), it verifies equation
\((1)\) for all generated deficit layers \(d\leq20\).

The code represents the monomial
\[
  q^j t^{M_{\ell,m}-d-j}
\]
by the pair \((j,M_{\ell,m}-d)\).  It creates three multisets:
\begin{itemize}
\item \emph{all}: one record for every generated path;
\item \emph{plus}: all positive string slots prescribed by the roots
      \(a_1=0\);
\item \emph{minus}: the correction slots coming from roots whose signed
      interval is oriented past its midpoint.
\end{itemize}
It then checks the exact multiset identity
\[
  \text{plus}=\text{all}+\text{minus}.
  \tag{2}
\]
This is a direct coefficient-by-coefficient verification of the conjectured
formula in the finite base range
\[
  0\leq d\leq20,
  \qquad
  \ell\leq\ell^*(m,20).
\]

Computation A has two logical uses.  First, it proves the theorem outright for
all \(\ell\leq\ell^*\).  Second, at \(\ell=\ell^*\), it says that the actual
paths occur at every area with exactly the multiplicity required by all the
proposed full string holsters.

In \texttt{code.py}, this is implemented by
\texttt{check\_lemma3}; the exposition-friendly alias is
\texttt{check\_computation\_a}.

\section{Computation B: no canonical string overshoots}

Computation B is Lemma 2 of Section 9; it is called ``computation 1'' in the
paper's appendix.  It is needed only at the critical length \(\ell^*\).

Let \(\lambda\in P_{\ell^*-1}^d\), and let \(h(\lambda)\) be the number of
parts of \(\lambda\).  The complete string indexed by \(\lambda\) is supposed
to end at area
\[
  R(\lambda)
  =
  M_{\ell^*,m}-d-h(\lambda).
\]
Computation B constructs the canonical endpoint
\[
  \lowest(f(g(\lambda)))
\]
and verifies
\[
  \area\bigl(\lowest(f(g(\lambda)))\bigr)
  \leq
  R(\lambda).
  \tag{3}
\]
Thus the canonical string may be too short, but it is never too long.  If it
passed its prescribed endpoint, the proof would have to shorten and regroup a
rigidly constructed string.  Inequality \((3)\) rules out that obstruction.

The program recognizes the roots by the condition \(a_1=0\), computes their
terminal paths with \texttt{lowest}, and recovers \(h(\lambda)\) through the
source statistic named \texttt{height}.  In \texttt{code.py}, the check is
\texttt{check\_lemma2}; the exposition-friendly alias is
\texttt{check\_computation\_b}.

\section{Why the two computations complete the boundary strings}

At \(\ell=\ell^*\), each canonical string fills a disjoint initial segment of
its prescribed holster.  Computation B proves that no initial segment extends
beyond the end of its holster.  Computation A proves that, at each area level,
the total number of actual paths equals the total number of holster slots.

Subtract the slots already occupied by the canonical strings.  At every area,
the number of unused paths is then exactly the number of unfilled slots.  Since
an extension tail is not required to follow \(\operatorname{right}\), the
unused paths can be assigned area by area to the open slots.  The result is a
partition of all paths into completed strings with the prescribed starts and
ends.  In the terminology of the paper, the layer is \emph{extendable}.

The construction is rigid in its canonical prefixes but flexible in its added
tails.  Correct multiplicities alone would prove the formula at the critical
length, but the canonical prefixes are essential for the symbolic propagation
to larger \(\ell\).

\section{Symbolic propagation}

Once every layer \(0\leq d\leq d^*\) is extendable at \(\ell^*\), Claim 16
of the paper proves symbolically that every such layer is extendable for all
\[
  \ell>\ell^*.
\]
The induction removes the largest parts from the partition indexing a root,
uses the extension already constructed one length lower, and prefixes the
appropriate new coordinate.  The proof checks symbolically that areas,
deficits, endpoints, and the unused-path partition all transform correctly.
No further computation is used after the critical length.

Consequently,
\[
  \text{Computation A below and at }\ell^*
  +
  \text{Computation B at }\ell^*
  +
  \text{symbolic propagation}
\]
proves the conjectured formula for every \(\ell\).

\section{Example: \(m=3\)}

Take \(m=3\) and \(d^*=20\).  The critical length is
\[
  \ell^*=8,
\]
because \(7\cdot3>20\), while \(6\cdot3\leq20\).

The proof specializes as follows.
\begin{enumerate}
\item \textbf{Finite base range.}  Computation A verifies the formula for
      every \(0\leq d\leq20\) and every \(\ell\leq8\).
\item \textbf{Canonical strings at the boundary.}  For each fitting partition
      of each \(d\), form \(f(g(\lambda))\) and apply
      \(\operatorname{right}\) until the path is unrightable.  The symbolic
      theory proves that these canonical strings are disjoint.
\item \textbf{No overshooting.}  Computation B checks that every canonical
      endpoint lies at or before its prescribed terminal area.
\item \textbf{Complete the tails.}  Computation A supplies the exact area
      multiplicities needed to allocate all leftover paths to the open slots.
      Hence every layer at \(\ell=8\) is extendable.
\item \textbf{Propagate.}  Claim 16 proves the formula for every \(\ell>8\)
      without further computation.
\end{enumerate}

\section{The finite run and its scope}

For \(d^*=20\), the appendix computation checks \(1\leq m\leq20\).  For
\(m>20\), the critical length is \(\ell^*=2\), and the paper handles the
remaining cases directly by symbolic arguments.

The recorded full run generated \(5{,}692{,}942\) source-relevant path
records.  Computation B found no failures.  Computation A checked 106
nontrivial \((m,\ell)\)-layers and found no failures.  The stable summary is
\[
\begin{array}{l}
\texttt{lemma2\_status: PASS},\quad
\texttt{lemma2\_failures: 0},\\
\texttt{lemma3\_status: PASS},\quad
\texttt{lemma3\_layers\_checked: 106},\quad
\texttt{lemma3\_failures: 0}.
\end{array}
\]

The readable file \texttt{code.py} closely follows the appendix routines and
uses only the Python standard library.  Its source defaults are
\texttt{--max-m 20} and \texttt{--dstar 20}.  Smaller checks can be run, for
example, with
\[
  \texttt{python code.py --max-m 3 --dstar 6}.
\]
The options \texttt{--check A} and \texttt{--check B} select the two
computations in the terminology of this note; \texttt{--lemma 3} and
\texttt{--lemma 2} are equivalent source-paper labels.

The theorem established here concerns the family \(r\equiv1\pmod s\) through
deficit \(20\).  The unrestricted formula for arbitrary coprime \((r,s)\),
and the same family in deficits above \(20\), remain conjectural.

\begin{thebibliography}{9}

\bibitem{Hawkes2024}
Graham Hawkes,
\emph{A conjectured formula for the rational \(q,t\)-Catalan polynomial},
Annals of Combinatorics \textbf{28} (2024), 749--795;
arXiv:2208.00577.

\end{thebibliography}

\end{document}
